\subsection{应用：表示半单性的判据}

定理 4. 设 \(\mathfrak{g}\) 是李代数，\(\mathfrak{r}\) 是其根基，\(\rho\) 是 \(\mathfrak{g},\mathfrak{g}' = \rho (\mathfrak{g})\) 的有限维表示，且 \(\mathfrak{r}' = \rho (\mathfrak{r})\)。则以下条件等价：

\begin{enumerate}
\def\labelenumi{(\alph{enumi})}
\item
  \(\rho\) 是半单的；
\item
  \(\mathfrak{g}'\) 是约化的，且其中心由半单自同态组成；
\item
  \(\mathfrak{r}'\) 由半单自同态组成；
\item
  \(\rho\) 在 \(\mathfrak{r}\) 上的限制是半单的。
\item
  \(\Rightarrow\) (b)：若 \(\rho\) 是半单的，则 \(\mathfrak{g}'\) 是约化的（命题 5）；由 1 和 \(\mathfrak{g}'\) 生成的结合代数是半单的（《代数》VIII 章，§ 5，第 1 号，命题 3），所以其中心是半单的（同上引文，§ 5，第 4 号，命题 12），从而该中心的元素都是半单的（同上引文，§ 9，第 1 号，命题 2）。
\item
  \(\Rightarrow\) (c)：若 \(\mathfrak{g}'\) 是约化的，则其中心等于其根基，即 \(\mathfrak{r}'\)，于是得到蕴含 (b) \(\Rightarrow\) (c)。
\item
  \(\Rightarrow\) (d)：假设 \(\mathfrak{r}'\) 由半单自同态组成。由于 \([\mathfrak{g}',\mathfrak{r}']\) 由幂零自同态组成（第 4 号，命题 6），有 \([\mathfrak{g}',\mathfrak{r}'] = \{0\}\)。然后，蕴含 (c) \(\Rightarrow\) (d) 由《代数》VIII 章§ 9，第 2 号，定理 1 得到。
\item
  \textgreater{} (a)：令 \(\mathfrak{s}\) 为 \(\mathfrak{g}\) 的幂零根，\(\rho'\) 为 \(\rho\) 在 \(\mathfrak{r}\) 上的限制。由于 \(\rho(\mathfrak{s})\) 的元素都是幂零的，\(\mathfrak{s}\) 包含于 \(\rho'\) 的最大幂零理想中。由于 \(\rho'\) 是半单的，\(\rho'(\mathfrak{s}) = \{0\}\)，而 \(\mathfrak{g}'\) 是约化的（命题 6 的推论），所以 \(\mathfrak{g}' = \mathfrak{a}' \times \mathfrak{r}'\)，其中 \(a'\) 是半单的（命题 5）。令 \(\mathbf{A}'\)（分别为 \(\mathbf{R}'\)）为由 1 和 \(a'\)（分别为 \(\mathfrak{r}'\)）生成的结合代数。它是半单的（《代数》VIII 章，§ 5，第 1 号，命题 3），因而 \(\mathbf{A}' \otimes_{\mathbb{K}} \mathbf{R}'\) 是半单的（同上引文，§ 7，第 6 号，定理 3 的推论 4），从而由 1 和 \(\mathfrak{g}'\) 生成的结合代数（它是 \(\mathbf{A}' \otimes_{\mathbb{K}} \mathbf{R}'\) 的一个商）是半单的，这就证明了 \(\rho\) 是半单的。
\end{enumerate}

推论 1. 设 g 是李代数，\(\rho\) 和 \(\rho'\) 是 g 的两个有限维半单表示。则 \(\rho\) 与 \(\rho'\) 的张量积是半单的。

令 \(\mathfrak{r}\) 为 \(\mathfrak{g}\) 的根基。对于 \(x \in \mathfrak{r}\)，\(\rho(x)\) 和 \(\rho'(x)\) 都是半单的（定理 4），因而 \(\rho(x) \otimes 1 + 1 \otimes \rho'(x)\) 是半单的（《代数》VIII 章§ 9，定理 1 的推论），所以 \(\rho\) 与 \(\rho'\) 的张量积是半单的（定理 4）。

推论 2. 设 \(\mathfrak{g}\) 是李代数，\(\rho\) 是 \(\mathfrak{g}\) 在有限维向量空间 \(V\) 上的半单表示，\(T\) 和 \(S\) 是 \(V\) 的张量代数与对称代数，\(\sigma_{T}\) 和 \(\sigma_{S}\) 是 \(\mathfrak{g}\) 在 \(T\) 和 \(S\) 上由 \(\rho\) 规范导出的表示。则 \(\sigma_{T}\) 和 \(\sigma_{S}\) 是半单的；更确切地说，它们是有限维单表示的直和。

令 \(T^{n}\) 为 T 中由 n 阶齐次张量组成的子空间。~该子空间在 \(\sigma_{T}\) 下稳定，而 \(\sigma_{T}\) 在 \(T^{n}\) 上定义的表示是半单的（推论 1）。因此关于 \(\sigma_{T}\) 的推论成立，从而关于 \(\sigma_{S}\) 的推论也成立，因为后者是 \(\sigma_{T}\) 的一个商表示。

推论 3. 设 g 是李代数，\(\rho\) 和 \(\rho'\) 是 g 在 M 和 M' 上的两个有限维半单表示。g 在 \(\mathcal{L}_{\mathrm{k}}(\mathrm{M},\mathrm{M}')\) 上的表示是由 \(\rho\) 和 \(\rho'\) 规范导出的，因而是半单的。

g-模 \(\mathcal{L}_{\mathbb{K}}(\mathbf{M},\mathbf{M}^{\prime})\) 可规范地与 g-模 \(\mathbf{M}^{*}\otimes_{\mathbb{K}}\mathbf{M}^{\prime}\) 等同（§ 3，第 3 号，命题 4），所以推论 3 由推论 1 得到。

推论 4. 设 \(\mathfrak{g}\) 是李代数，a 是 \(\mathfrak{g}\) 的一个理想，\(\rho\) 是 \(\mathfrak{g}\) 的半单表示。

\begin{enumerate}
\def\labelenumi{(\alph{enumi})}
\item
  \(\rho'\) 为 \(\rho\) 在 a 上的限制是半单的。
\item
  若 \(\rho\) 是单的，则 \(\rho'\) 是彼此同构的单表示之和。
\end{enumerate}

将 \(\rho\) 对其核取商后，可以假设 \(\rho\) 是忠实的。于是 \(\mathfrak{g}\) 是约化的。令 \(\mathfrak{g} = \mathfrak{g}_1 \times \mathfrak{g}_2\)，其中 \(\mathfrak{g}_1\) 是 \(\mathfrak{g}\) 的中心，\(\mathfrak{g}_2\) 是半单的。于是 \(\mathfrak{a} = \mathfrak{a}_1 \times \mathfrak{a}_2\)，其中 \(\mathfrak{a}_1 \subset \mathfrak{g}_1\)，\(\mathfrak{a}_2 \subset \mathfrak{g}_2\)，且 \(\mathfrak{a}_1\) 是 \(\mathfrak{a}\) 的中心。\(\rho(\mathfrak{g}_1)\) 的元素，特别是 \(\rho(\mathfrak{a}_1)\) 的元素，都是半单的（定理 4），因而 \(\rho\) 是半单的（定理 4）。于是得到 (a)。断言 (b) 由 (a) 得到，应用§3，第 1 号，命题 1 的推论。

